\subsection{Life-Cycle model 29: Exploiting Monotonicity for Faster Solutions (Divide-and-Conquer)}
There is no new model here. Instead we just add a single line of code to (almost) any of the Life-Cycle Models, namely
\\ \textit{vfoptions.divideandconquer=1}
\\ \noindent and as a result, it solves more than twice as fast! In fact the model also uses less GPU memory when solving, meaning it is not only faster, but that you can solve bigger models on a given GPU!\footnote{Because of how GPUs parallelize there is an exception, namely for very small models it may actually be slower.} Most of the models in this Intro to Life-Cycle Models actually already have this activated.

The codes for Life-Cycle Model 29 just solve the same model as Life-Cycle Model 9, twice. The first time uses the defaults and stores the resulting value and policy functions as $V1$ and $Policy1$. The second time sets \textit{vfoptions.divideandconquer=1}, turning on divide-and-conquer, and stores the results as $V2$ and $Policy2$. Each call to the value function iteration is timed using $tic/toc$, so you can see for yourself how much faster the second call is. After both solutions are computed we check that $max(abs(V1(:)-V2(:)))$ and $max(abs(Policy1(:)-Policy2(:)))$ are both zero, which confirms that as long as 'conditional monotonicity' is satisfied, using divide-and-conquer gives the exact same solution (just faster, and less GPU memory used).

Divide-and-conquer works by exploiting the monotonicity of the policy function, that is, it is exploits that $aprime(a)$ is a monontone increasing function (that next periods choice of endogenous state is a monotone increasing choice of this periods endogenous state). To be more precise, it requires a weaker conditional montonicity, so we require that $aprime(a; d,z)$ is a monotone increasing function (that next periods choice of endogenous state is a monotone increasing choice of this periods endogenous state, conditional on values for d and z). This is a property that almost all Economic models satisfy, although there are some expections.

In practice, you would want to set \textit{vfoptions.divideandconquer=1} in every single one of the models we have seen (it is possible that it would be slower in the really small models, the exact criterion of 'really small' depends on the specs of your GPU). These models all satisfy the monotonicity assumption, and we can solve them faster by exploiting this property.

Divide-and-conquer in VFI Toolkit works in two 'layers'. \textit{vfoptions.level1n=round(sqrt(n\_{}a));} is the (default) number of points (in the grid on this period endogenous state) used in the first layer. If you increase/decrease this number you can make the codes go faster still (the best number depends heavily on the model, so you just have to play around with it). The speed gains of changing are modest compared to the gains from turning on divide-and-conquer, so in most cases you can just not declare \textit{vfoptions.level1n} and leave it at default value unless you really do need the speed.

 On rare occasions there are some non-monotone models used in Economics, and so it is important that you think about whether your model is monotone before you activate \textit{vfoptions.divideandconquer=1}. If you are unsure if your model satisfies the monotonicity assumption, you have two options: (i) sit down and think carefully through the equations, (ii) run the model twice with \textit{vfoptions.divideandconquer=0} and \textit{vfoptions.divideandconquer=1} and check that both give identical answers (which they only will if the model is monotone; if the answers differ the model is non-monotone and only the \textit{vfoptions.divideandconquer=0} answer is correct). Obviously the first of these is the more through and better, but the second is very easy and convenient (albeit not completely certain as it only illustrates when conditional monotonicity fails, but cannot show that it holds).

For a more comprehensive explanation of what divide-and-conquer is and what exactly the technical monotonicity assumption involves, see:
\href{http://discourse.vfitoolkit.com/t/divide-and-conquer-solve-one-endogenous-state-models-faster/265}{http://discourse.vfitoolkit.com/t/divide-and-conquer-solve-one-endogenous-state-models-faster/265}


\subsection{Life-Cycle model 30: Linear Interpolation of $aprime$ (Grid Interpolation Layer)}
There is no new model here either.  Instead we just add four lines of code to (almost) any of the Life-Cycle Models, namely
\\ \textit{vfoptions.gridinterplayer=1}
\\ \textit{vfoptions.ngridinterp=20}
\\ \textit{simoptions.gridinterplayer=vfoptions.gridinterplayer;}
\\ \textit{simoptions.ngridinterp=vfoptions.ngridinterp;}
\\ \noindent and the toolkit will let $aprime$ take values between the grid points on $a\_{}grid$ using linear interpolation. The first line turns on the grid interpolation layer; the second sets the number of evenly-spaced points placed between each pair of consecutive points of $a\_{}grid$ (you can change $20$ to any integer you like). Unlike with divide-and-conquer (Life-Cycle Model 29), turning on the grid interpolation layer does change the answer: $V$ differs and even the size of $Policy$ changes (there is an extra entry storing the position of $aprime$ on the interpolation layer, between two consecutive $a\_{}grid$ points). For any given $n\_{}a$ the grid interpolation layer is marginally slower but much more accurate, because $aprime$ is no longer restricted to lie on $a\_{}grid$. In practice, turning on grid interpolation layer allows you to use a smaller $n\_{}a$ than would otherwise be required, and so it ends up allowing your codes to be both faster and more accurate, and even reduces the GPU memory use.

The codes for Life-Cycle Model 30 just solve the same model as Life-Cycle Model 9, three times. The first time uses the defaults and stores the resulting value and policy functions as $V1$ and $Policy1$. The second time sets \textit{vfoptions.gridinterplayer=1} and \textit{vfoptions.ngridinterp=20}, turning on the grid interpolation layer, and stores the results as $V2$ and $Policy2$. We then check $size(Policy1)$ and $size(Policy2)$ to see that the grid interpolation layer adds one extra entry to $Policy$, and we check $max(abs(V1(:)-V2(:)))$ to see that the two value functions differ. The third times uses both divide-and-conquer and grid interplation layer and is just to emphasize that both can be used together.

We then compute the agent distribution twice. Because the grid interpolation layer changes the size and interpretation of $Policy$, we also need to tell the stationary distribution command about it by setting \textit{simoptions.gridinterplayer=vfoptions.gridinterplayer} and \textit{simoptions.ngridinterp=vfoptions.ngridinterp}. We then check $max(abs(StationaryDist1(:)-StationaryDist2(:)))$, the maximum absolute difference between the two stationary distributions: this is small but not zero.

If you set $n\_{}a$ to a really large value, then the values of $aprime$ that the grid interpolation layer can pick will be very close to $a\_{}grid$ points anyway, so the two solutions become essentially identical. You can try setting $n\_{}a$ to a large value (e.g., $1001$ or $2001$) at the top of the codes and rerunning; you will see that the differences in $V$ and $StationaryDist$ between the two solutions get smaller and smaller as $n\_{}a$ gets larger.

For a more comprehensive explanation of what the grid interpolation layer is and what it does, see:
\href{http://discourse.vfitoolkit.com/t/grid-interpolation-layer-vfoptions-gridinterplayer-1/401}{http://discourse.vfitoolkit.com/t/grid-interpolation-layer-vfoptions-gridinterplayer-1/401}
